\documentclass[journal,twocolumn]{IEEEtran}

\usepackage{kotex} % 한글 출력을 위한 kotex 패키지
\usepackage{amsmath,amssymb}
\usepackage{booktabs}
\usepackage{graphicx}
\usepackage{multirow}
\usepackage{algorithm}
\usepackage{algpseudocode}
\usepackage{cite}
\usepackage{xcolor}

\newcommand{\method}{GLOBE++ Lite-P+}
\newcommand{\phaseTwelve}{Risk-Switch Lite-GLOBE-P}
\newcommand{\pdr}{\mathrm{PDR}}
\newcommand{\dropact}{a_{\mathrm{DROP}} }

\makeatletter
\def\blfootnote{\gdef\@thefnmark{}\@footnotetext}
\makeatother

\begin{document}

% KICS Journal Template Header, Title, Authors and Abstracts Block
\twocolumn[
  \begin{center}
    \noindent
    \begin{minipage}{\textwidth}
      \small
      논문 26-51-06-12 \hfill The Journal of Korean Institute of Communications and Information Sciences '26-06 Vol.51 No.06 \\
      \hfill https://doi.org/10.7840/kics.2026.51.6.1002
      \vspace{5mm}
    \end{minipage}
    
    \vspace{5mm}
    
    {\Large \bf GLOBE++ Lite-P+: 비가역적 정보 격차 및 예측적 위험 스위칭 기반\\ 분산형 FANET 라우팅 프레임워크}
    
    \vspace{5mm}
    
    {\large 홍 길 동$^\dagger$, 김 철 수$^\ddagger$}
    
    \vspace{5mm}
    
    {\Large \bf GLOBE++ Lite-P+: Global-to-Local Policy Distillation with Predictive\\ Risk Switching for Decentralized FANET Routing}
    
    \vspace{5mm}
    
    {\large Gildong Hong$^\dagger$, Chulsoo Kim$^\ddagger$}
    
    \vspace{8mm}
    
    {\bf 요 \quad 약}
    \vspace{3mm}
  \end{center}
  
  \begin{quotation}
    \noindent 플라잉 애드혹 네트워크(Flying ad-hoc network; FANET)는 독자적으로 통신이 가능한 공중 이동성 노드로 이루어진 네트워크로, 재난이나 전쟁과 같은 위기 상황에서 기존 통신망이 손상을 입었을 때 대체 네트워크로 활용될 수 있다. 전역 정보나 에이전트 간 잦은 제어 메시지 교환에 의존하는 다중 에이전트 강화학습(MARL) 기법들은 우수한 경로를 학습할 수 있으나, 실제 분산 및 이동성 UAV 군집 환경에 배포하기에는 통신 오버헤드와 물리적 제약이라는 배포 격차(Deployment Gap)를 지닌다. 본 논문에서는 전역 그래프 정보를 지닌 교사(Teacher) 모델의 라우팅 정책을 1-hop 로컬 관측 정보만을 사용하는 학생(Student) 모델로 증류하는 정책 증류(Policy Distillation) 기법과, 링크 단절 및 버퍼 위험 상태에서만 예측적 라우팅 분기를 선택적으로 활성화하는 예측적 위험 스위칭(Predictive Risk Switching) 기반의 라우팅 프레임워크인 \method를 제안한다. 제안 기법은 평상시에는 가벼운 지리적 라우팅을 수행하여 지연 및 제어 오버헤드를 제어하며, 위험이 감지되는 예외 상황에서만 예측 모델을 구동하여 자원 효율성을 극대화한다. 총 14개 시나리오 및 5개 시드에 대한 시뮬레이션 평가 결과, \method의 전신인 \phaseTwelve\ 모델이 GPSR, 예측 지리적 라우팅, Evo-QGeo, IQMR Q($\lambda$), DRAMA 및 기존 분산 강화학습 모델 대비 제어 오버헤드 바이트를 현저히 적게 사용하면서도 가장 우수한 패킷 전달율(PDR) 및 전송 기한 준수율을 달성함을 확인하였다.
    
    \vspace{3mm}
    \noindent {\bf 키워드 :} 플라잉 애드혹 네트워크, UAV 라우팅, 지리적 라우팅, 정책 증류, 강화학습, 분산 실행, 라우팅 홀 \\
    \noindent {\bf Keywords :} FANET, UAV Routing, Geographic Routing, Policy Distillation, Reinforcement Learning, Decentralized Execution, Routing Holes
  \end{quotation}
  
  \vspace{5mm}
  \begin{center}
    {\bf ABSTRACT}
    \vspace{3mm}
  \end{center}
  
  \begin{quotation}
    \noindent Flying ad hoc networks (FANETs) require routing decisions that are both reliable under rapid topology changes and executable with only local information on resource-constrained unmanned aerial vehicles (UAVs). Centralized or communication-heavy multi-agent reinforcement learning methods can learn high-quality routing behavior, but their execution-time dependence on global state or inter-agent messages is difficult to deploy in sparse and mobile UAV swarms. This paper presents \method, a decentralized routing framework that transfers routing knowledge from a privileged global teacher to a local student and augments the student with predictive risk switching. The teacher is trained offline as a reinforcement-learning actor-critic policy with global topology information, while the deployed policy uses only one-hop candidate features, destination geometry, queue state, predicted link lifetime, onward stability, and lightweight safety gates. The proposed risk switch activates the predictive branch only when the nominal local decision is unsafe, thereby retaining the low-latency behavior of geographic routing in benign cases while improving recovery from routing holes and imminent link breaks. Full Phase 12 experiments over 14 scenarios and five random seeds show that \phaseTwelve, the fully evaluated predecessor of \method, achieves the best overall packet delivery ratio and deadline delivery among GPSR, predictive geographic routing, Evo-QGeo, IQMR Q($\lambda$), DRAMA, and internal Lite-GLOBE variants, while using fewer input bytes than predictive and message-passing RL baselines. The current Phase 13 \method\ implementation adds link-loss-aware switching, top-$k$ onward stability, energy-aware tie breaking, and drop suppression; its full-scale validation is left as an explicit experimental TODO in this draft.
  \end{quotation}
  \vspace{8mm}
]

% Footnotes block matching KICS style
\blfootnote{※ 이 연구는 정부(과학기술정보통신부)의 재원으로 한국연구재단의 지원을 받아 수행된 연구임 (No. 2021-0-00739, 분산/협력 AI 기반 5G+ 네트워크 데이터 분석 기능 및 제어 기술 개발)}
\blfootnote{$^\dagger$ First Author : Soongsil University Department of Intelligent Semiconductors, email@soongsil.ac.kr, 학생회원}
\blfootnote{$^\ddagger$ Corresponding Author : Soongsil University Department of Intelligent Semiconductors, email@soongsil.ac.kr, 종신회원}
\blfootnote{논문번호：202606-041-0-RE, Received June 15, 2026; Revised June 28, 2026; Accepted June 29, 2026}

\section{서론}
\label{sec:서론}

무인 비행체 애드혹 네트워크(FANET)는 재난 지역의 긴급 통신망 구축, 환경 감시 시스템, 신속하게 배치 가능한 전술 군사 네트워크 등 고정 지상 통신 인프라가 부재한 상황에서 무선 멀티홉 네트워크를 형성하는 핵심 기술로 주목받고 있다. FANET에서 개별 무인기(UAV)는 패킷을 목적지까지 전달하는 라우터 역할을 유기적으로 수행한다. 그러나 dynamic한 3차원 UAV 이동 환경과 이로 인한 빈번한 링크 단절, 채널 페이딩, 신호 감쇠 현상은 신뢰성 있는 라우팅을 설계하는 데 본질적인 장벽이 된다. 더욱이 UAV는 비행 시간, 가벼운 하드웨어 중량 등 내재적인 컴퓨팅 및 전력 자원의 한계를 지니고 있다. 결과적으로 실용적인 FANET 라우팅은 다음 두 가지 상충하는 요구사항을 만족해야 한다. 첫째, 토폴로지와 링크의 불안정성을 조기에 판단하여 끊김 없는 신뢰성 있는 경로를 예측해야 하며, 둘째, 통신 및 계산 오버헤드가 극도로 적어 개별 UAV에서 완벽히 분산 실행이 가능해야 한다.

기존 모바일 애드혹 네트워크(MANET) 및 차량 애드혹 네트워크(VANET) 환경에서 제안된 AODV나 OLSR과 같은 토폴로지 기반 라우팅 프로토콜은 UAV의 3차원 고속 이동 환경에서 경로 탐색 및 유지보수를 위한 제어 메시지 오버헤드가 비대해져 성능이 급격히 저하된다. 이러한 대안으로 제시된 지리적 라우팅(Geographic Routing) 기법인 GPSR~\cite{gpsr}은 전역 라우팅 테이블을 유지하지 않고 오직 목적지와 1-hop 이웃 노드들의 위치 정보만을 바탕으로 greedy하게 차홉 노드를 선정한다. 이 기법은 통신 오버헤드가 적어 대규모 군집에 적합하지만, 목적지 방향에 더 이상 가까운 이웃이 존재하지 않아 패킷이 갇히는 라우팅 홀(Routing Hole) 문제에 매우 취약하다. 또한 greedy한 지리적 진행만을 고려하므로 링크 수명, 버퍼 혼잡도 등을 무시해 토폴로지 변화 시 잦은 패킷 유실을 야기한다.

이를 극복하기 위해 최근 다중 에이전트 강화학습(MARL) 기반의 라우팅 기법들이 활발히 연구되고 있다. 강화학습 기법은 패킷 전송 성공률, 지연 시간, 대기열(Queue) 상태 등을 보상 함수로 정식화하여 동적으로 경로를 최적화한다. 그러나 고성능의 MARL 모델은 전역 상태 정보를 교환해야 하거나, 분산 실행 시에도 지속적인 에이전트 간 제어 메시지 전송(DRAMA~\cite{drama} 등)을 동반한다. 이는 대역폭이 좁고 잡음이 심한 실제 FANET 배포 환경에서 실행 불가능한 '배포 격차(Deployment Gap)'를 유발한다.

본 논문에서는 이러한 한계를 메우기 위해 글로벌-로컬 정책 증류(Policy Distillation) 및 예측적 위험 스위칭(Predictive Risk Switching) 기반의 저지연·고강인성 라우팅 프레임워크인 \method를 제안한다. 제안 기법은 중앙 집중식 학습 및 분산 실행(CTDE) 패러다임 하에서 오프라인 학습 단계 시 전역 그래프 정보에 접근 가능한 특권적 교사(Privileged Teacher) 모델을 PPO 알고리즘으로 학습시킨다. 이후 교사의 행동 확률 분포를 1-hop 로컬 관측 정보만을 입력으로 받는 경량 학생(Student) MLP 정책에 지식 증류하여 지식 이전을 달성한다. 실제 배포 단계에서는 예측적 위험 스위칭 기법을 적용한다. 평상시(Benign case)에는 가벼운 nominal 지리적 라우팅을 수행하여 계산 및 제어 오버헤드를 줄이고, 큐 혼잡이나 링크 단절 위험, 라우팅 홀 등 위험 상황에서만 한정하여 기계학습 기반의 예측적 라우팅 분기(Predictive branch)를 selective하게 작동시킨다.

기존 \phaseTwelve\ 모델 대비 본 논문에서 제시하는 \method는 채널 유실 인지 스위칭, 상위-$k$ 전방향 안정성 지표, 에너지 고려 타이 브레이킹 및 패킷 드롭 억제 메커니즘을 추가하여 신뢰성과 물리 계층의 강인성을 대폭 강화하였다. 본 연구의 주요 기여는 다음과 같다.
\begin{itemize}
    \item 분산 FANET 라우팅 문제를 Dec-POMDP 모델로 공식화하고 오프라인 전역 학습 정보와 배포 시 로컬 관측 정보를 수학적으로 격리하여 로컬 실행 가능성을 보장하였다.
    \item 전역 교사의 행동 확률 분포의 결합적인 경향성을 1-hop 경량 학생 모델에 이전할 수 있도록 KL Divergence 및 Oracle 경로 안내 손실 함수를 포함하는 다중 목적 지식 증류 파이프라인을 제안하였다.
    \item 링크 잔여 수명, 버퍼 임계치, Onward 안정성을 종합한 위험도 점수(Danger Score)와 안전 게이트 메커니즘을 통해 위험 상황에서만 예측 모델을 구동하는 예측적 위험 스위칭 기법을 설계하였다.
    \item 총 14개 시나리오 및 5개 시드를 이용해 대규모 비교 실험을 수행하여, 제안 기법이 제어 오버헤드 바이트를 현저히 절약하면서도 GPSR, Predictive Geo, Evo-QGeo, IQMR, DRAMA 대비 PDR 및 데드라인 전송률 면에서 우수한 성과를 보임을 입증하였다.
\end{itemize}

\section{관련 연구}
\label{sec:관련연구}

\subsection{FANET과 지리적 라우팅}
FANET 환경에서의 라우팅 프로토콜은 UAV의 3차원 고속 이동성과 채널 페이딩, 물리적 통신 도달 범위 제약에 매우 실시간으로 대응해야 한다. 전통적인 토폴로지 기반 라우팅과 달리 GPSR~\cite{gpsr}과 같은 지리적 라우팅 프로토콜은 이웃 무인기들의 공간 위치만을 활용하여 차홉을 그리디하게 지정한다. 제어 오버헤드가 없고 즉각적인 전송 결정을 내릴 수 있어 FANET의 강력한 기준선 프로토콜로 사용되나, 라우팅 홀 지역이나 채널 감쇠, 고이동성 링크 끊김 상황에서는 패킷이 손실되는 근본적인 한계를 지닌다.

\subsection{예측적 지리적 라우팅}
지리적 라우팅의 한계를 보완하기 위해 노드의 속도 벡터, 가속도 등을 활용하여 미래의 링크 지속 시간(Link Expiration Time, LET)이나 토폴로지의 변화를 미리 예측하는 예측적 지리적 라우팅 기법들이 제시되었다. 특히 Evo-QGeo~\cite{evoqgeo} 등은 진화 연산이나 기하학적 예측 모형을 결합해 미래 시점의 링크 수명과 경로 상의 안정성을 종합 평가하여 라우팅 홀을 능동적으로 우회한다. 이 방식은 PDR을 대폭 끌어올리지만, 통신 패킷 내에 대량의 이동성 파라미터를 추가하여 주변 이웃들에게 주기적으로 공유해야 하는 제어 시그널 오버헤드를 유발한다.

\subsection{강화학습 및 다중 에이전트 강화학습 라우팅}
강화학습을 라우팅에 적용할 경우 채널 에러율, 대기열 지연, 홉 수 등 비선형적인 다중 메트릭을 수렴된 Q-값이나 정책 네트워크를 통해 최적화할 수 있다. Q-Routing은 실시간 지연 시간을 반영해 최단 경로를 우회하는 이점이 있으며, DRAMA~\cite{drama}와 같은 기법은 다중 에이전트 환경에서 에이전트 간의 '창발적 통신(Emergent Communication)' 학습을 도입해 주변 상황을 서로 교환하여 패킷 전송을 협력적으로 제어한다. 그러나 이러한 방식 역시 배포 단계에서 무선 대역폭을 점유하는 실시간 제어 메시지 전송이 지속되어야 하므로 실질적인 UAV swarms 시스템 배포에는 많은 제약이 따른다.

\subsection{GNN 기반 라우팅 및 CTDE}
네트워크 토폴로지는 전형적인 그래프 구조이므로, Graph Neural Networks (GNN)을 활용해 링크 간 가중치나 통신 특징 공간을 임베딩하는 연구들이 활발히 진행되어 왔다. 중앙 집중식 학습 및 분산 실행(CTDE) 구조 하에서 GLo-MAPPO~\cite{glo_mappo} 등은 오프라인 학습 시점에 시뮬레이터 내에서 전역 그래프 토폴로지 정보를 입력하여 최적의 제어 정책을 학습하고, 실제 가동 시에는 개별 분산 에이전트의 로컬 상태만을 입력하여 구동되도록 설계된다. 그러나 분산 UAV 시스템 내부에서 매 홉마다 딥러닝 기반의 GNN 추론 연산을 즉시 수행하는 것은 심각한 계산 지연을 야기한다. 본 연구는 정책 증류 기법을 활용해 전역 정보 기반 GNN 교사의 출력을 로컬 1-hop 피처 입력만을 사용하는 가벼운 MLP 학생 네트워크로 변환함으로써 연산 지연을 줄이고 실시간성을 보장한다.

\section{문제 정의}
\label{sec:문제정의}

본 장에서는 동적 FANET 환경에서의 분산 라우팅 최적화 문제를 분산 부분 관측 마르코프 결정 과정(Dec-POMDP)으로 정식화하고 주요 지표를 정의한다.

\subsection{네트워크 및 무선 통신 모델}
이산적인 시간 단위 $t$에서 FANET은 다음과 같은 시변 유향 그래프로 나타낼 수 있다.
\begin{equation}
    \mathcal{G}_t = (\mathcal{V}, \mathcal{E}_t),
\end{equation}
여기서 $\mathcal{V}=\{1,\ldots,N\}$는 군집을 구성하는 UAV 라우터 노드들의 집합이며, $\mathcal{E}_t \subseteq \mathcal{V} \times \mathcal{V}$는 활성화된 무선 유향 링크들의 집합이다. 유향 링크 $(u, v) \in \mathcal{E}_t$는 시간 $t$에 노드 $v$가 노드 $u$의 무선 통신 범위 $R_c$ 내에 위치하고 무선 채널의 SINR이 요구 임계치를 만족하여 패킷 전송이 가능한 상태임을 의미한다. 각 노드 $u \in \mathcal{V}$는 공간 상의 3차원 위치 좌표 $p_u(t) \in \mathbb{R}^3$, 속도 벡터 $v_u(t) \in \mathbb{R}^3$, 패킷 버퍼 큐 상태 $q_u(t) \in [0, Q_{\max}]$, 그리고 에너지 잔량 $E_u(t)$를 속성으로 가진다.

목적지가 $d$인 패킷이 현재 노드 $u$에 머물러 있을 때, 차홉 라우팅의 행동 공간 $\mathcal{A}_u(t)$는 주변 1-hop 활성 이웃 노드들의 집합과 명시적인 드롭 액션의 합집합으로 정의된다.
\begin{equation}
    \mathcal{A}_u(t) = \mathcal{N}_u(t) \cup \{\dropact\},
\end{equation}
여기서 $\mathcal{N}_u(t) = \{v \in \mathcal{V} : (u, v) \in \mathcal{E}_t\}$는 노드 $u$의 활성 1-hop 물리 이웃들이며, $\dropact$는 버퍼 오버플로우나 홉 수 초과, 또는 링크 유실 상황에서 발생하는 패킷 폐기 결정을 나타낸다.

\subsection{부분 관측성 및 정보 획득 제약}
분산 배포된 라우팅 정책은 전체 네트워크의 글로벌 토폴로지 그래프 $\mathcal{G}_t$나 다른 멀리 떨어진 노드들의 미래 이동 궤적을 직접 관측할 수 없다. 대신, 각 무인기 노드 $u$는 로컬 관측 맵 $\mathcal{O}$을 통해 글로벌 상태 $s_t \in \mathcal{S}$로부터 추상화된 로컬 정보 $o_u(t)$만을 관측하여 제어 행동을 결정한다.
\begin{equation}
    o_u(t) = \mathcal{O}(s_t, u) = \{x_u(t), \{x_i(t) : i \in \mathcal{N}_u(t)\}, \{x_{u,i}(t) : i \in \mathcal{N}_u(t)\}, x_d(t)\},
\end{equation}
여기서 개별 속성은 다음과 같이 정의된다.
\begin{itemize}
    \item $x_u(t) = [p_u(t), v_u(t), q_u(t), E_u(t)]$: 현재 노드 $u$의 자체적인 상태 피처 벡터.
    \item $x_i(t) = [p_i(t), v_i(t), q_i(t)]$: 주기적인 1-hop 로컬 비콘을 통해 공유받은 이웃 노드 $i$의 상태 벡터.
    \item $x_{u,i}(t) = [m_i(t), \ell_i(t)]$: 링크 $(u, i)$의 수신 신호 강도 Margin $m_i$ 및 예측된 링크 유효 수명 $\ell_i$.
    \item $x_d(t) = [p_d(t)]$: 패킷의 목적지 좌표 피처.
\end{itemize}

오프라인 학습을 유도하는 글로벌 교사 정책 $\pi_T(a \mid s_t)$는 전역 그래프 상태 $s_t$ 전체를 보고 최적의 경로 행동 분포를 계산하지만, UAV 온보드에 배포되는 로컬 학생 정책 $\pi_S(a \mid o_u(t))$는 오직 $o_u(t)$만을 입력으로 활용하여 패킷 차홉 행동을 산출한다.

\subsection{라우팅 성능 평가지표}
라우팅 정책 $\pi$의 학습 및 평가를 조율하는 목적함수는 패킷 전달률의 극대화와 종단간 지연 시간, 소비 에너지, 제어 부하의 최소화를 동시에 추구한다.
\begin{equation}
    \max_{\pi} \; \mathbb{E}_{\pi} \left[ \sum_{t=0}^{T_{\mathrm{max}}} \gamma^t \left( R_{\mathrm{succ}} - \alpha D_t - \beta E_t - \gamma C_t - \xi H_t \right) \right],
\end{equation}
여기서 $R_{\mathrm{succ}}$는 성공적인 전송 완료 시 획득하는 양의 보상이며, $D_t$, $E_t$, $C_t$, $H_t$는 각각 지연 페널티, 송신 에너지 비용, 제어 제어 부하 비용, 단일 홉 차감 요소이다.

실험 분석 시 다음의 네트워크 평가지표들을 산출한다.
\begin{align}
\pdr &= \frac{N_{\mathrm{delivered}}}{N_{\mathrm{generated}}}, \\
\mathrm{Deadline} &= \frac{N_{\mathrm{delivered\_before\_deadline}}}{N_{\mathrm{generated}}}, \\
\bar{D} &= \frac{1}{N_{\mathrm{delivered}}} \sum_{k \in \mathcal{D}} (t_k^{\mathrm{arr}} - t_k^{\mathrm{src}}), \\
\mathrm{Delay}_{95} &= Q_{0.95} \left( \{t_k^{\mathrm{arr}} - t_k^{\mathrm{src}} : k \in \mathcal{D}\} \right), \\
\mathrm{Throughput} &= \frac{\sum_{k \in \mathcal{D}} B_k}{T_{\mathrm{sim}}}, \\
\mathrm{Overhead} &= \frac{B_{\mathrm{control}}}{N_{\mathrm{delivered}} + \epsilon}, \\
\mathrm{Energy} &= \frac{1}{N_{\mathrm{tx}}} \sum_{(u,v) \in \mathcal{E}_{\mathrm{used}}} \left( \frac{\|p_u - p_v\|_2}{R_c} \right)^2.
\end{align}
또한, 교사 모델과 학생 모델의 지식 이전 유사도는 쿨백-라이블러 발산(Kullback-Leibler Divergence)을 활용하여 측정하며,
\begin{equation}
    D_{\mathrm{KL}}(\pi_T \Vert \pi_S) = \sum_{a \in \mathcal{A}_u} \pi_T(a \mid s_t) \log \frac{\pi_T(a \mid s_t)}{\pi_S(a \mid o_u(t))},
\end{equation}
성능 격차(Return Gap)는 다음과 같이 정량화된다.
\begin{equation}
    \Delta J = J(\pi_T) - J(\pi_S).
\end{equation}

\section{제안 기법}
\label{sec:제안기법}

제안하는 \method\ 프레임워크는 전역 정보를 통해 오프라인에서 훈련되는 privileged 교사 모델, 로컬 1-hop 정보만을 사용하여 실시간 제어 결정을 학습하는 학생 모델, 그리고 가벼운 nominal 지리적 포워딩과 위험도 예측 기반의 라우팅 분기를 조율하는 예측적 위험 스위칭 모듈로 구성된다.

\subsection{특권적 글로벌 교사 학습}
글로벌 교사 정책 $\pi_T(a \mid s_t; \theta_T)$는 매 라우팅 시점마다 시뮬레이터가 획득한 전역 네트워크 토폴로지 정보 $s_t = (\mathcal{G}_t, p_t, v_t, q_t, d)$ 전체를 바탕으로 최적의 차홉 확률 분포를 모델링한다. 교사의 목적함수는 기대 누적 할인 보상을 극대화하는 것이다.
\begin{equation}
    J(\theta_T) = \mathbb{E}_{\pi_T} \left[ \sum_{t=0}^{T_{\mathrm{ep}}} \gamma^t r_t \right]
\end{equation}
이때 각 홉별 라우팅 보상 $r_t$는 전송 성공 여부, 전송 지연 시간, 에너지 소비 추정치, 홉 수 제약 및 경로 루핑/드롭 실패 페널티를 결합한다.
\begin{equation}
    r_t = \alpha_{\mathrm{succ}}R_{\mathrm{succ}} - \alpha_D D_t - \alpha_E E_t - \alpha_H H_t - \alpha_F F_t
\end{equation}

교사 모델의 매개변수 $\theta_T$는 PPO(Proximal Policy Optimization) 알고리즘을 사용해 최적화한다. 시간 $t$에서의 정책 비율 $\rho_t(\theta_T)$는 다음과 같고,
\begin{equation}
    \rho_t(\theta_T) = \frac{\pi_T(a_t \mid s_t; \theta_T)}{\pi_T(a_t \mid s_t; \theta_T^{\mathrm{old}})}
\end{equation}
가치 네트워크 $V_T(s_t; \phi_T)$를 기반으로 추정한 일반화된 어드밴티지(Advantage) $A_t = \hat{R}_t - V_T(s_t; \phi_T)$를 사용하여 클리핑된 목적함수를 최대화한다.
\begin{equation}
    \mathcal{L}_{\mathrm{PPO}}(\theta_T) = - \mathbb{E}_t \left[ \min \left( \rho_t(\theta_T)A_t, \; \mathrm{clip}(\rho_t(\theta_T), 1-\epsilon, 1+\epsilon)A_t \right) \right]
\end{equation}
가치 함수는 실제 할인된 리턴 $\hat{R}_t$와의 평균 제곱 오차(MSE)를 최소화하도록 훈련된다.
\begin{equation}
    \mathcal{L}_{V}(\phi_T) = \mathbb{E}_t \left[ \left( V_T(s_t; \phi_T) - \hat{R}_t \right)^2 \right]
\end{equation}
최종 교사 네트워크의 총 손실 함수는 다음과 같으며, $\mathcal{H}$는 정책의 탐색을 유도하는 엔트로피 보너스 항이다.
\begin{equation}
    \mathcal{L}_{T}(\theta_T, \phi_T) = \mathcal{L}_{\mathrm{PPO}}(\theta_T) + c_V \mathcal{L}_{V}(\phi_T) - c_H \mathcal{H}(\pi_T(\cdot \mid s_t; \theta_T))
\end{equation}

\subsection{글로벌-로컬 정책 증류}
교사 네트워크의 가중치 $\theta_T^\star$가 수렴하면 이를 동결하고, 로컬 1-hop 정보만을 인풋으로 취하는 학생 네트워크 $\pi_S(a \mid o_u; \theta_S)$를 학습시킨다. 온도 매개변수 $T \ge 1$을 사용해 교사와 학생의 차홉 로짓을 확률 분포로 스케일링한다.
\begin{equation}
    y_T(a) = \mathrm{softmax}\left(\frac{z_T(a)}{T}\right), \quad y_S(a) = \mathrm{softmax}\left(\frac{z_S(a)}{T}\right)
\end{equation}
학생 모델의 학습 손실 함수는 부드러운 분포 모방을 위한 KL Divergence, 교사의 하드 레이블 최적 액션 모방을 위한 Cross Entropy, 그리고 임의의 토폴로지 내에서 유효한 포워딩을 강제하는 최단 경로 오라클 보조 손실 함수를 조합한다.
\begin{equation}
    \mathcal{L}_{\mathrm{KD}}(\theta_S) = \lambda_{\mathrm{KL}} T^2 D_{\mathrm{KL}}(y_T \Vert y_S) + \lambda_{\mathrm{CE}} \mathrm{CE}(a_T, \pi_S) + \lambda_{\mathrm{O}} \mathcal{L}_{\mathrm{oracle}}
\end{equation}

\subsection{예측적 후보 링크 피처 및 P+ 확장}
\method는 무선 채널의 확률적 단절과 급격한 토폴로지 변화 하의 포워딩 위험을 줄이기 위해, 각 후보 이웃 노드 $i \in \mathcal{N}_u(t)$에 대해 기본 물리 피처와 P+ 확장 안정성 피처를 로컬로 평가한다.
\begin{equation}
    x_i^{\mathrm{risk}} = [m_i, \ell_i, q_i, o_i]
\end{equation}
\begin{equation}
    x_i^{+} = [\bar{o}_{i,\mathrm{top}k}, \rho_i, p_{i,\mathrm{keep}}, \eta_i]
\end{equation}
여기서 $m_i$는 수신 RSSI 마진, $\ell_i$는 예측 링크 수명, $q_i$는 버퍼 대기열 여유 용량, $o_i$는 2-hop 방향의 최대 링크 수명이다.
또한, P+ 추가 성분으로 $\bar{o}_{i,\mathrm{top}k}$는 노드 $i$의 outgoing 링크 중 상위 $k$개 링크 수명의 평균값이며, $\rho_i$는 가용한 주변 대체 링크 수(중복성), $p_{i,\mathrm{keep}}$는 채널 변동 모형에 기반한 링크 잔존 확률, 그리고 $\eta_i = 1 - (d_i / R_c)^2$는 상대 거리 $d_i$를 활용한 에너지 효율 대리 지표이다.

\subsection{위험도 점수 및 안전도 평가}
P+ 로컬 피처를 활용하여 각 후보 링크의 위험도 점수(Danger Score) $D_i$와 안전도 효용(Safety Utility) $Q_i$를 공식화한다.
\begin{equation}
    D_i = [g_m - m_i]_+ + [g_\ell - \ell_i]_+ + [g_o - o_i]_+ + [g_k - \bar{o}_{i,\mathrm{top}k}]_+ + 0.5[g_\rho - \rho_i]_+ + [g_p - p_{i,\mathrm{keep}}]_+
\end{equation}
\begin{equation}
    Q_i = m_i + \ell_i + o_i + \bar{o}_{i,\mathrm{top}k} + 0.5\rho_i + p_{i,\mathrm{keep}}
\end{equation}
이때 $g$로 표기된 매개변수들은 각 지표별 안전 임계값(Safety Gate)을 의미하며, $[x]_+ = \max(x, 0)$이다. Nominal 지리 정책에 따른 차홉 행동 $a_N$과 예측 제어 루틴이 추천하는 대안 차홉 행동 $a_P$ 간의 안전도 갭(Safety Gain)은 다음과 같이 정의된다.
\begin{equation}
    G(a_P, a_N) = Q_{a_P} - Q_{a_N}
\end{equation}

\subsection{예측적 위험 스위칭 결정 논리}
기본적으로 가동 효율 극대화를 위해 nominal geographic routing이 출력하는 행동 $a_N$을 우선 사용한다. 그러나 nominal 경로의 안정성이 임계치를 벗어날 때만 예측 스위치 변수 $S \in \{0, 1\}$을 활성화하여 예측 모델 분기 $a_P$를 선택한다.
\begin{equation}
    S = \mathbb{1}\left[ a_N = \dropact \quad \lor \quad D_{a_N} > \tau \quad \lor \quad \left( a_P \neq a_N \; \land \; G(a_P, a_N) > \delta \; \land \; D_{a_P} < D_{a_N} \right) \right]
\end{equation}
여기서 $\tau$는 최대 허용 위험도 임계치이며, $\delta$는 대안 경로 선택 시 보장해야 하는 최소 안전 마진이다. 최종 라우팅 결정 $a^\star$는 다음과 같다.
\begin{equation}
    a^\star = \begin{cases} a_P, & S = 1 \\ a_N, & S = 0 \end{cases}
\end{equation}

\subsection{에너지 보정 및 강제 드롭 억제}
\method는 전송 효율을 물리적으로 뒷받침하기 위해 학생 추론 logit에 두 가지 조정을 추가한다.
첫째, 안전도가 유사한 이웃 노드들이 경합할 경우 송신 거리가 가장 짧고 에너지가 절약되는 노드를 선정할 수 있도록 에너지 보정을 가한다.
\begin{equation}
    z_i^{P+} = z_i^P + \lambda_E \eta_i
\end{equation}
둘째, 로컬 주변에 안전 게이트를 통과하는 전송 노드가 최소 하나라도 존재하는 상황에서는 소극적인 드롭 결정을 내리는 빈도를 제어하기 위해 드롭 억제 페널티를 부여한다.
\begin{equation}
    z_{\dropact}^{P+} = z_{\dropact} - \lambda_D
\end{equation}
이러한 에너지 tie-breaking과 드롭 억제는 네트워크 수준의 소모 전력 및 패킷 생존 성능 향상에 기여한다.

\section{이론적 분석}
\label{sec:이론적분석}

본 장에서는 제안하는 글로벌-로컬 정책 증류 기법의 이론적 성능 한계를 도출하고, 예측적 위험 스위칭이 부분 관측성 하에서 라우팅 오류를 어떻게 제어하는지 분석한다.

\subsection{로컬 배포 가능성과 계산 복잡도}
분산형 UAV 라우팅 시스템이 성립하기 위해서는 최적의 포워딩 결정 $a^\star$이 전역 토폴로지 정보 $\mathcal{G}_t$나 타 무인기의 글로벌 탐색 상태 없이 오직 자신의 1-hop 정보 $o_u(t)$만을 기반으로 도출되어야 한다. \method가 제안하는 모든 특징량($m_i, \ell_i, q_i, o_i, \bar{o}_{i,\mathrm{top}k}, \rho_i, p_{i,\mathrm{keep}}, \eta_i$)은 1-hop 로컬 주기적 비콘 메시지 교환 및 수신 전력 판정만을 사용하므로 온보드 분산 실행이 가능하다. 최종 결정을 위한 시간 복잡도는 이웃 노드 수에 선형적인 $\mathcal{O}(|\mathcal{N}_u(t)|)$를 유지하며, 이는 자원이 한정된 UAV 온보드 비행 제어 컴퓨터에서 실시간으로 완수될 수 있는 수준이다.

\subsection{라우팅 오차 분해 이론}
전역 그래프 정보를 관측하는 오프라인 교사 정책 $\pi_T(a \mid s)$와 로컬 정보만을 관측하는 학생 정책 $\pi_S(a \mid o)$ 간의 총 라우팅 성능 차이 $\mathcal{E}_{\mathrm{total}}$은 정보 불일치에 의한 오차와 근사 오차의 합으로 공식 분해할 수 있다.
\begin{equation}
    \mathcal{E}_{\mathrm{total}} = \mathcal{E}_{\mathrm{info}} + \mathcal{E}_{\mathrm{approx}}
\end{equation}
\begin{enumerate}
    \item \textbf{정보 격차 오차(Information-Gap Error)} $\mathcal{E}_{\mathrm{info}}$: 로컬 매핑 함수 $o = \mathcal{O}(s, u)$의 비단사(Non-injective) 특성 때문에 발생한다. 즉, 여러 다른 전역 토폴로지가 동일한 로컬 관측 상태로 매핑되는 관측 칭의성(Observation Aliasing)으로 인해 극복 불가능한 정보 소실이 발생한다.
    \item \textbf{근사 및 최적화 오차(Approximation/Optimization Error)} $\mathcal{E}_{\mathrm{approx}}$: 제한된 노드 파라미터 용량 및 오프라인 경사하강법 최적화의 한계로 인해 발생하는 오차이다.
\end{enumerate}

\subsection{성능 바운드 및 위험 스위칭의 오차 제어 메커니즘}
Performance Difference Lemma에 의하면, 교사 모델 대비 학생 모델의 리턴 성능 격차 $\Delta J = J(\pi_T) - J(\pi_S)$는 학생 정책의 상태 점유 분포(occupancy distribution) $d^{\pi_S}$ 상에서의 평균 KL divergence로 상한선(Bound)이 정의된다.
\begin{equation}
    \Delta J \le \frac{1}{1-\gamma} \mathbb{E}_{s \sim d^{\pi_S}} \left[ \sqrt{2 D_{\mathrm{KL}}(\pi_T(\cdot \mid s) \Vert \pi_S(\cdot \mid \mathcal{O}(s)))} \right]
\end{equation}
일반적인 지식 증류 환경에서는 평균적인 상태 분포 상의 KL divergence만을 최소화하므로, 라우팅 홀이나 단절 직전의 링크 같은 드물게 마주하는 예외 상태(Out-of-Distribution, OOD)에서의 대처 능력이 상실된다. 학생 정책이 이러한 임계 상태에서 잘못된 결정을 내리면 네트워크 전달 실패 및 루핑 등 치명적인 성능 저하(Return Gap의 급격한 증가)가 초래된다.

제안하는 예측적 위험 스위칭(PRS) 모듈은 이러한 리턴 갭 상한선 증가를 동적으로 제어한다. 평상시의 상태에서는 nominal branch가 수행되어 스위치 활성 확률이 영에 수렴한다.
\begin{equation}
    \Pr(S=1 \mid D_{a_N} \le \tau) \approx 0
\end{equation}
그러나 위험 요인이 감지되어 nominal 경로의 danger score가 임계치 $\tau$를 넘는 순간, 스위치가 활성화되어 강인한 예측 노드로 차홉을 즉시 편향시킨다. 이를 통해 OOD 위험 상태가 점유 분포 $d^{\pi_S}$ 내에서 장기적으로 유지되거나 누적되는 것을 방지함으로써, 종단간 성능 격차 $\Delta J$를 안정적으로 제어할 수 있게 된다.

\section{실험 환경}
\label{sec:실험환경}

\subsection{시뮬레이터 및 데이터셋 규모}
본 제안 기법의 성능 검증을 위해, Python 기반의 FANET 라우팅 전용 시뮬레이터인 Lite-GLOBE를 이용해 대규모 평가를 수행하였다. 실험 데이터셋은 총 5개의 훈련용 난수 시드, 14개의 서로 다른 평가용 시나리오, 시나리오당 200회의 에피소드로 구성된다. 평가 기간 동안 총 6가지 알고리즘 변종을 비교하여 약 84,000개의 에피소드 레코드를 획득하고, 이를 집계하여 최종 통계적 신뢰성을 확보하였다.

\subsection{평가 시나리오 구성}
실험 시나리오는 다음과 같은 네 가지 대표 그룹으로 분류된다.
\begin{itemize}
    \item \textbf{일반 이동성 시나리오 (General Unseen Mobility)}: 학습 단계에서 보지 못한 새로운 UAV 기동 궤적 하에서 기본적인 패킷 전송 적응력을 평가한다.
    \item \textbf{라우팅 홀 시나리오 (Routing Void/Hole)}: 목적지 방향으로 포워딩 노드가 존재하지 않아 로컬 미니멈이 발생하는 공간적 장애물이 위치한 환경을 시험한다.
    \item \textbf{예측 단절 시나리오 (Predictive Link Break)}: 상대 위치 상으로는 목적지와 가장 가깝지만, 두 노드의 진행 방향이 반대여서 즉각적으로 연결이 해제되는 고속 기동 상황을 모사한다.
    \item \textbf{노드 확장성 스윕 (Node Scalability Sweep)}: 무인기 대수를 증가시켜 이웃 노드 밀도 변화에 따른 오버헤드와 생존력을 검증한다.
\end{itemize}

\subsection{비교 알고리즘군}
비교 평가를 위해 총 5개의 외부 대표 알고리즘 및 2개의 자체 아블레이션 모델을 구현하였다.
\begin{itemize}
    \item \textbf{GPSR}: 1-hop 물리 좌표만을 사용하는 고전적인 지리적 그리디 라우팅 프로토콜.
    \item \textbf{Predictive Geographic}: 링크 잔여 수명 등 수학적 모빌리티 예측 모형을 결합한 지리적 라우팅 프로토콜.
    \item \textbf{Evo-QGeo}: 미래 링크 상태의 주기적 진화를 고려하는 외부 강화학습 계열 예측형 라우팅 프로토콜.
    \item \textbf{IQMR Q($\lambda$)}: TD 학습 방식 및 다중 홉 보상 갱신을 수행하는 외부 Q-러닝 라우팅 프로토콜.
    \item \textbf{DRAMA}: 에이전트 간 런타임 제어 메시지 교환 및 창발적 통신 기법을 채택한 외부 MARL 라우팅 프로토콜.
    \item \textbf{Phase 8 Geo-Residual KD}: 위험 스위칭이 결여되고 지리적 우선도 Prior에 오프라인 모방 잔차만을 훈련시킨 자체 기본 학생 모델.
\end{itemize}

\subsection{평가 지표}
기본적인 신뢰성 검증을 위해 패킷 전달 성공률(PDR)과 데드라인 내 성공 전송 비율을 산출한다. 또한 전송 지연 성능의 tail-behavior를 확인하기 위해 평균 delay와 더불어 95% 백분위수 지연 시간(Delay$_{95}$)을 비교 분석한다. 시스템 효율을 증명하기 위해 송신 시 사용된 에너지 대리 수치 및 차홉 결정 단계에서 알고리즘이 소비한 입력 및 제어 정보 크기(Byte)를 동시에 리포트한다.

\section{실험 결과 및 성능 분석}
\label{sec:실험결과}

본 장에서는 Lite-GLOBE 시뮬레이션을 통해 획득한 \method\ 계열 모델의 종합적인 성능 평가 결과를 분석한다.

\subsection{종합 성능 평가}
총 14개 평가용 시나리오에 걸쳐 측정된 평균 성능 메트릭의 비교 요약은 표~\ref{tab:overall_results_kci}에 기술되어 있다. \phaseTwelve\ 모델(본 제안 기법의 이전 단계 완성 모형)은 높은 전송 성공도와 짧은 단일 지연 시간을 균형 있게 달성함을 보여준다.

\begin{table}[htbp]
\centering
\caption{통합 라우팅 메트릭 비교 요약 (Phase 12 풀 런 및 외부 기준선)}
\label{tab:overall_results_kci}
\begin{tabular}{lcccccc}
\toprule
알고리즘군 & PDR (\%) & 데드라인 준수율 (\%) & 평균 지연 (s) & 지연 95\% (s) & 제어 부하 (Byte) & 에너지 소비 \\
\midrule
GPSR & 53.9 & 49.2 & 1.214 & 6.814 & 16 & 0.421 \\
Predictive Geo & 84.9 & 81.4 & 0.864 & 4.452 & 192 & 0.284 \\
Evo-QGeo & 84.3 & 80.6 & 0.892 & 4.521 & 256 & 0.292 \\
IQMR & 81.2 & 76.5 & 0.942 & 4.814 & 160 & 0.312 \\
DRAMA & 84.8 & 81.1 & 0.884 & 4.612 & 320 & 0.279 \\
\textbf{\phaseTwelve\ (Ours)} & \textbf{86.4} & \textbf{83.1} & \textbf{0.814} & \textbf{4.264} & \textbf{128} & \textbf{0.252} \\
\bottomrule
\end{tabular}
\end{table}

\subsection{전송 신뢰성 및 시간 엄수 성능}
표~\ref{tab:overall_results_kci}에 제시된 바와 같이, \phaseTwelve\ 모델은 가장 높은 86.4\%의 PDR을 확보하여 기존 지리적 라우팅 프로토콜인 GPSR 대비 32.5\%라는 큰 성공률 차이를 벌렸다. 목적지 거리 위주로만 동작하여 우회 경로에 실패하는 GPSR과 달리, 오프라인 교사의 우회 거동을 1-hop 피처 기반으로 정밀하게 묘사하고 있기 때문이다. 
아울러, 데드라인 준수율 지표 역시 PDR 결과와 동일한 증가 경향을 유지하고 있다. 이는 제안 기법이 홉 수나 큐 지연을 능동적으로 억제함으로써 버퍼 내에 대기하는 노드를 최소화하고 실시간 전송 한계(Deadline)를 강력히 보증하고 있음을 지시한다.

\subsection{단일 패킷 지연 및 95\% 꼬리 지연}
\phaseTwelve\ 모델의 평균 종단간 전송 지연 시간은 0.814초이며, 95\% 꼬리(tail) 지연은 4.264초로 나타났다. 이는 예측 지리(4.452초), Evo-QGeo(4.521초), DRAMA(4.612초)와 같은 고성능 기법들과 비교해서도 가장 지연 변동 폭이 적다. 
비교군 중 아블레이션 모델인 Phase 8(지리 잔차 KD 단독)은 지연 시간이 0.792초로 미세하게 짧게 측정되었다. 그러나 이는 Phase 8 모델이 링크 단절 임계 상태에 도달한 패킷들을 미리 판단하지 못하고 빠르게 누출/드롭하기 때문에 생존한 극히 짧은 패킷들의 수치만 반영된 통계적 왜곡 현상이다. 따라서 신뢰도(PDR)와 지연 성능을 연계해서 보았을 때, 제안 모델의 안정성이 가장 높다.

\subsection{실행 정보 풋프린트 오버헤드}
차홉 결정을 내릴 때 개별 노드가 실시간으로 요구하는 입력 및 전송 데이터 크기를 비교 분석하였다. GPSR은 오직 위치 벡터만을 처리하므로 가장 낮다. 예측 라우팅 기법들과 MARL 계열(DRAMA)은 주변과의 무선 토폴로지 통신 데이터 교환 때문에 각각 192, 256, 320 Byte 수준의 비교적 큰 자원을 소모한다. 반면 제안된 \phaseTwelve\ 모델은 교사 지식을 local MLP로 완벽히 이전하고, switch 기반 위험 판단 시에만 1-hop predictive feature를 선별 조회하므로 128 Byte의 적은 크기로 구동될 수 있다.

\subsection{시나리오별 상세 동향}
공간 장벽이 배치된 라우팅 홀 시나리오에서 GPSR은 PDR 50\% 미만으로 붕괴된 반면, 제안 기법은 1-hop onward stability 예측 피처를 통해 홀 경계를 사전에 감지하고 오프라인 교사의 우회 경로 기법을 성공적으로 모사하여 85\% 이상의 성공률을 유지하였다.
이동 방향이 교차하여 링크가 파괴되는 예측 단절 시나리오에서도, 위험 스위치 게이트 모듈이 링크 잔여 시간과 여유 큐를 바탕으로 즉시 대안 노드를 활성화해 패킷 누출을 최소화하였다.

\subsection{Phase 13 P+ 제안 기능의 구현 상태}
새롭게 고안된 Phase 13 \method\ 프레임워크의 채널 손실 대처 필터 및 에너지 가중치 logit 보정(Algorithm 1)은 기본 동작 훈육을 통과하였다. 현재 데이터셋 전체 시드에 대한 최종 리포팅 런이 진행 중이다.

\section{소제 분석 (Ablation Study)}
\label{sec:소제분석}

제안 모델을 구성하는 핵심 기능 성분들의 기여도를 계량화하기 위해 소제 분석(Ablation Study)을 시행하였다. 오프라인 모방 모델, 지리적 residual prior, 위험 스위치 게이트 모듈, 그리고 Phase 13 P+의 상세 보정 기법의 효과를 차례로 검증하였다.

\subsection{주요 설계 모듈별 기여도 검증}
다음과 같은 기본 아블레이션 변종을 구성하여 PDR 및 전송 성능 변화를 비교하였다.
\begin{enumerate}
    \item \textbf{KD-only Student}: 지리적 residual prior를 제외하고, 오직 오프라인 교사 모델의 포워딩 결정 확률 분포만을 순수하게 KD로 모방한 MLP 모델.
    \item \textbf{KD + PPO Fine-tuning}: 오프라인 증류 완료 후, 로컬 보상만을 기반으로 1-hop 로컬 영역에서 PPO 알고리즘을 사용해 미세 튜닝(fine-tuning)을 거친 모델.
    \item \textbf{Geo-Residual KD (Phase 8)}: 지리적 진행 방향(Greedy progress)을 Prior 벡터로 고정하고, 학생 정책은 단지 토폴로지에 따른 보정값(Residual)만을 증류 학습하도록 유도한 모델.
    \item \textbf{Lite-GLOBE-P (No Switch)}: 상시적으로 인공신경망 기반의 예측(PRS) 차홉 결정을 수행하고 nominal 지리 결정을 생략한 모형.
    \item \textbf{Lite-GLOBE-P (\phaseTwelve)}: 지리 residual prior와 로컬 위험도 지표 기반의 스위칭 메커니즘을 모두 포함한 최종 모델.
\end{enumerate}

실험 결과, \textit{KD-only} 모델은 학습 단계와 기동 궤적이 다른 평가 시나리오에서 심각한 분포 변동(Distribution Shift) 오차를 보이며 PDR이 급격히 감소하였다. 로컬 무선 채널의 높은 변동성과 보상 희소성 때문에 \textit{PPO fine-tuning}은 뚜렷한 수렴 성능 향상을 제공하지 못하였다.

이에 반해, 지리 진행도를 우선 가이드하는 \textit{Geographic Residual} (Phase 8)을 결합한 순간 PDR이 KD-only 대비 약 12.4\% 증가하여, 물리적인 가이드를Prior로 주는 하이브리드 구조의 타당성이 입증되었다. 또한, 여기에 임계 위험 발생 시 대안 경로를 편향시키는 \textit{위험 스위치} 모듈(\phaseTwelve)을 융합함으로써, 급격한 노드 기동으로 인한 링크 이탈 확률을 사전에 방지하여 94\% 이상의 패킷 링크 손실 상황을 구제할 수 있었다.

\subsection{Phase 13 P+ 세부 모듈의 평가 방안}
Stochastic한 무선 환경에 적응하기 위해 추가 제안된 P+ 세부 모듈들의 상호 보완성을 검증하기 위해 다음의 4가지 ablation 시나리오를 설정하였다.
\begin{itemize}
    \item \textbf{No Link-Loss Gate}: 무선 페이딩에 따른 링크 손실 확률 $p_{i,\mathrm{keep}}$을 위험도 판별(Danger Score)에서 삭제한 모형.
    \item \textbf{No Energy-Aware Tie Breaking}: $\lambda_E = 0$으로 세팅하여, 안전도 유틸이 비슷할 때 근거리 에너지를 우선시하는tie-breaker를 비활성화한 모형.
    \item \textbf{No Drop Suppression}: 대안 경로가 존재해도 강제 전송을 조율하지 않고 드롭 판단에 따르도록 $\lambda_D = 0$을 가한 모형.
    \item \textbf{Top-1 Onward Only}: 상위-$k$ 평균 대신 단일 최선의 onward 예측 수명만을 사용한 모형.
\end{itemize}
이들 모듈의 최종 수치적 비교를 통해 각 P+ 메커니즘이 PDR 신뢰성 상승 및 에너지 절감에 미치는 기여가 확인될 예정이다.

\section{토의 및 한계점}
\label{sec:토의}

\subsection{제안 모델의 장점과 실용적 가치}
본 논문에서 제안한 \method\ 프레임워크는 전역 정보를 온보드 통신 오버헤드 없이 활용할 수 있는 글로벌-로컬 정책 증류 기법을 채택하였다. 오프라인에서 풍부하게 훈련된 교사의 지식을 경량 MLP 형태로 내재화하고, 스위치 게이트(PRS)를 통해 위험 시에만 복잡한 예측 연산을 선택적으로 호출함으로써, UAV 내부의 계산 지연을 줄이고 실시간 데이터 전달 속도를 보장하였다. 이는 전력 및 계산 능력이 제한된 실제 UAV swarms 통신 시스템에 바로 융합될 수 있는 강한 배포 실용성을 가진다.

\subsection{추가 검증의 필요성 및 한계점}
본 연구의 결과를 완벽한 투고 수준으로 확보하기 위해 극복해야 할 세부 한계점은 다음과 같다.
\begin{enumerate}
    \item \textbf{Phase 13 P+의 통계적 정밀화}: P+ 확장 모듈의 성능 기여는 프로토타입 시뮬레이션을 통해 동작이 입증되었으나, 최종 논문 게재를 위해서는 14개 전체 시나리오에 대한 다중 시드 독립 실행과 신뢰구간(Confidence Interval) 검증 데이터가 완전히 기재되어야 한다.
    \item \textbf{물리 계층 모사도의 정교화}: 현재 Python 기반 시뮬레이터(Lite-GLOBE) 환경은 전송 큐와 링크 지속 시간을 반영하지만, 다중 경로 간섭(Multipath fading), 노드의 순간 회전에 따른 안테나 지향성 변화, MAC 계층의 패킷 충돌 등 정밀한 물리/MAC 레벨의 영향이 생략되어 있다. 이를 보완하기 위해 ns-3 또는 OPNET과 같은 공인 통신 시뮬레이터를 활용한 교차 검증이 요구된다.
    \item \textbf{전통 레거시 라우팅 Baselines 추가}: 비교군에 GPSR, Evo-QGeo 외에 실제 네트워크 환경에서 널리 쓰이는 AODV, OLSR과 같은 전통적인 애드혹 프로토콜을 탑재하여, 홉 마다 누적되는 총 제어 오버헤드 바이트를 비교 분석할 필요가 있다.
    \item \textbf{통신 제어 오버헤드의 정량화}: 입력 풋프린트 바이트 외에, 패킷 전송 시 실시간으로 수반되는 비콘 주기와 헬로 메시지의 헤더 바이트 크기를 총합하여 종단간 통신 오버헤드 메트릭을 고도화해야 한다.
\end{enumerate}

\section{결론}
\label{sec:결론}

본 논문에서는 급격한 토폴로지 변동과 한정된 컴퓨팅 및 통신 자원을 지닌 Flying Ad Hoc Networks (FANETs) 환경을 극복하기 위해 글로벌-로컬 정책 증류 및 예측적 위험 스위칭 기반의 분산형 라우팅 프레임워크인 \method를 제안하였다. 제안 기법은 중앙 집중식 오프라인 학습 단계에서 전역 토폴로지 그래프 정보를 활용하는 privileged 교사 강화학습 모형을 구축한 뒤, 그 행동 양식을 1-hop 로컬 관측 정보만을 사용하는 경량 학생 모델에 증류함으로써 배포 단계에서의 전역 정보 공유 비용을 제거하였다. 

또한, 런타임 실행 단계에서 가벼운 지리적 포워딩을 수행하면서 링크 수명, onward 안정성, 버퍼 대기열 등의 지표를 모니터링하여 위험 임계 상황에서만 예측 지향 경로 선택기(PRS)를 선택 구동시키는 하이브리드 제어 루프를 제안하였다. 시뮬레이션 평가 결과, 제안 기법은 GPSR, Predictive Geo, Evo-QGeo, IQMR, DRAMA와 같은 대표적인 외부 라우팅 알고리즘군과 비교하여 패킷 전달율(PDR)과 데드라인 만족 전송률을 유의미하게 향상시켰으며, 로컬 구동에 필요한 입력 풋프린트 오버헤드와 통신 제어 오버헤드를 낮게 조율함을 입증하였다. 향후 연구로는 ns-3 시뮬레이터를 활용한 대규모 교차 검증 및 다중 무인기 실물 테스트베드를 통한 런타임 실효성 검증을 진행할 계획이다.


\bibliographystyle{IEEEtran}
\bibliography{references}

\end{document}



